%%%%%%%%%%%%%%%%%%%%%%%%%%%%%%%%%%%%%%%%%%%%%%%%%%%%%%%%%%%%%%%%%%%%%
%% This is a "template" model document for submission to the
%% American Chemical Society (ACS).
%%
%% The guidance here contains information about how you may wish to
%% modify it to match the requirements of various ACS journals. The
%% ACS do *not* typeset accepted articles using LaTeX, so there is
%% no specific class required.
%%
%% This template deliberately does *not* seek to reproduce
%% the layout of the typeset journal: this is explicitly not
%% required by the ACS for LaTeX submissions.
%%
%% Please report any issues with the template at
%% https://github.com/josephwright/acs-template/issues
%%
%% Released under the Creative Commons 0 license
%% https://creativecommons.org/public-domain/cc0/
%% 
%% Copyight (c) 2025 Joseph Wright
%%%%%%%%%%%%%%%%%%%%%%%%%%%%%%%%%%%%%%%%%%%%%%%%%%%%%%%%%%%%%%%%%%%%%
\documentclass[letterpaper]{article}

%%%%%%%%%%%%%%%%%%%%%%%%%%%%%%%%%%%%%%%%%%%%%%%%%%%%%%%%%%%%%%%%%%%%%
%% Font setup - delete if you are using LuaLaTeX
%%%%%%%%%%%%%%%%%%%%%%%%%%%%%%%%%%%%%%%%%%%%%%%%%%%%%%%%%%%%%%%%%%%%%
\usepackage[T1]{fontenc}

%%%%%%%%%%%%%%%%%%%%%%%%%%%%%%%%%%%%%%%%%%%%%%%%%%%%%%%%%%%%%%%%%%%%%
%% Adjust the margins and allow for line spacing
%%%%%%%%%%%%%%%%%%%%%%%%%%%%%%%%%%%%%%%%%%%%%%%%%%%%%%%%%%%%%%%%%%%%%
\usepackage{geometry}
\geometry{margin = 1in}
\usepackage{setspace}
\usepackage{bm}
%%%%%%%%%%%%%%%%%%%%%%%%%%%%%%%%%%%%%%%%%%%%%%%%%%%%%%%%%%%%%%%%%%%%%
%% Reference support
%%
%% The recommended method for producing the reference section is
%% to use biblatex. If you wish to use a classical BibTeX
%% approach, this is easiest to achieve using the achemso package.
%% In that case, you should remove the biblatex lines.
%%%%%%%%%%%%%%%%%%%%%%%%%%%%%%%%%%%%%%%%%%%%%%%%%%%%%%%%%%%%%%%%%%%%%
% You can adjust the printing of DOI, article title, etc. using
% package options, e.g. "doi = true"; see the biblatex manual for
% details of adjusting the number of authors printed, e.g.
% "maxnames = 15" to print no more than 15 authors.
% \usepackage[style = chem-acs]{biblatex}
% \addbibresource{pqe.bib}
% If you are using classical BibTeX, remove the above lines and 
% uncomment:
\usepackage{achemso}

%%%%%%%%%%%%%%%%%%%%%%%%%%%%%%%%%%%%%%%%%%%%%%%%%%%%%%%%%%%%%%%%%%%%%
%% Graphic inclusion and scheme and chart support3....
%%%%%%%%%%%%%%%%%%%%%%%%%%%%%%%%%%%%%%%%%%%%%%%%%%%%%%%%%%%%%%%%%%%%%
\usepackage{graphicx}
\usepackage{float}
\newfloat{scheme}{htbp}{los}
\floatname{scheme}{Scheme}
\floatname{chart}{Chart}
\newfloat{graph}{htbp}{loh}

%%%%%%%%%%%%%%%%%%%%%%%%%%%%%%%%%%%%%%%%%%%%%%%%%%%%%%%%%%%%%%%%%%%%%
%% Common support packages
%%%%%%%%%%%%%%%%%%%%%%%%%%%%%%%%%%%%%%%%%%%%%%%%%%%%%%%%%%%%%%%%%%%%%
\usepackage{chemformula} % Formulas using \ch{}
% or
\usepackage[version = 4]{mhchem} % Formulas using \ce{}

%%%%%%%%%%%%%%%%%%%%%%%%%%%%%%%%%%%%%%%%%%%%%%%%%%%%%%%%%%%%%%%%%%%%%
%% Many journals require that sections are unnumbered: this 
%% is activated here
%%%%%%%%%%%%%%%%%%%%%%%%%%%%%%%%%%%%%%%%%%%%%%%%%%%%%%%%%%%%%%%%%%%%%
\setcounter{secnumdepth}{-1}

%%%%%%%%%%%%%%%%%%%%%%%%%%%%%%%%%%%%%%%%%%%%%%%%%%%%%%%%%%%%%%%%%%%%%
%% Place any additional macros here.  Please use \newcommand* where
%% possible, and avoid layout-changing macros (which are not used
%% when typesetting).
%%%%%%%%%%%%%%%%%%%%%%%%%%%%%%%%%%%%%%%%%%%%%%%%%%%%%%%%%%%%%%%%%%%%%
\newcommand*\mycommand[1]{\texttt{\emph{#1}}}

%%%%%%%%%%%%%%%%%%%%%%%%%%%%%%%%%%%%%%%%%%%%%%%%%%%%%%%%%%%%%%%%%%%%%
%% Author and title data:
%% the authblk package is currently the simplest way to provide this
%%%%%%%%%%%%%%%%%%%%%%%%%%%%%%%%%%%%%%%%%%%%%%%%%%%%%%%%%%%%%%%%%%%%%
\usepackage{authblk}
\author[1]{Patryk Kozlowski}
\affil[1]{Department of Chemistry and Chemical Biology, Harvard University, Cambridge, MA 02138, USA}


\title{Theoretical Spectroscopy via Many-Body Perturbation Theory}

\begin{document}

\maketitle


%%%%%%%%%%%%%%%%%%%%%%%%%%%%%%%%%%%%%%%%%%%%%%%%%%%%%%%%%%%%%%%%%%%%%
%% Start the main part of the manuscript here.
%%%%%%%%%%%%%%%%%%%%%%%%%%%%%%%%%%%%%%%%%%%%%%%%%%%%%%%%%%%%%%%%%%%%%
\abstract{
}\section{Introduction}
\label{sec:intro}
Core spectroscopy is notoriously difficult to simulate. The sudden creation of a hole couples to the plasmonic modes of the system, forming an electronic polaron. As a result, the observed spectrum consists of a sharp quasiparticle peak accompanied by satellite structures that arise from the coupling of the hole to the plasmons of the system. This spectrum is described by the spectral function, which is proportional to the imaginary part of the retarded Green's function.

Many-body perturbation theory, the state-of-the-art approach for excited-state calculations, is commonly used to obtain the spectral function. Specifically, it constructs the Green's function through a perturbative expansion in the screened Coulomb interaction $W$. Its lowest-order form, known as $G_0W_0$, has been empirically successful in describing the position and lifetime of the quasiparticle peak, but it fails to describe the satellite features. One might hope to remedy this by performing the expansion self-consistently; however, self-consistent $GW$ does not necessarily yield a better description of the satellites than $G_0W_0$.
% TODO: cite self-consistent GW satellite studies

Another option is coupled-cluster theory, which has become the gold standard For computing ground-state properties in molecular quantum chemistry. Excited-state properties, such as spectral functions, can be obtained from its equation-of-motion extension (EOM-CC). However, the  steep computational scaling of EOM-CC limits its application to small supercells.

A viable alternative is the cumulant expansion, which resums the perturbation series, thereby incorporating a class of vertex corrections. As a result, it qualitatively reproduces the core spectrum, generating a sequence of plasmon satellites. The cumulant expansion is typically performed on top of a $G_0W_0$ calculation, an approach we refer to as $GW$+C. Consequently, $GW$+C inherits the starting-point dependence of $G_0W_0$. Furthermore, studies of finite supercells of the electron gas have shown that $G_0W_0$ deviates significantly from numerically exact results.
% TODO: cite finite-supercell benchmarks

A remedy lies in the observation that core spectroscopy gives rise to an electronic polaron, in which the suddenly created hole couples to the bosonic degrees of freedom of the system. In this work, we exploit this observation by writing down a physically motivated Hamiltonian for the electronic polaron (albeit deriving a controlled approximation to the \emph{ab initio} problem) and solving for its Green's function. We are aided in this endeavor by the rapid development, in recent years, of techniques that provide near-exact solutions of polaron problems at reasonable computational cost.
% \subsubsection{TODO: Why do we study the electron gas at rs = 4; It is a good model for alkali metals, and this density is also interesting because of the outside correlation contribution. Such that we are no longer in the perturbative regime
% }

In Sec.~\ref{sec:theory}, we introduce the theories used to tackle this problem, beginning with the state-of-the-art $GW$+C approach, followed by our quasiboson Hamiltonian and the techniques we use to solve it, namely the self-consistent cumulant expansion and delocalized coupled-cluster theory. In Sec.~\ref{sec:results}, we analyze the results obtained from these approaches. Finally, in Sec.~\ref{sec:conclusions}, we summarize our findings and discuss future directions.

\section{Theory}
\label{sec:theory}
We compute spectral functions of the homogeneous electron gas at a Wigner--Seitz radius of $r_s = 4$. The electron gas is a bellwether system for excited-state methods, whose ultimate goal is to predict the spectral functions of real materials. This density is of practical interest because it coincides with the valence electron density of the simple alkali metals. It is also of theoretical interest because, at $r_s = 4$, the electron gas lies outside the perturbative regime, where there is a notable lack of existing  methods that can accurately compute excited states.

The theories investigated below make use of the screened Coulomb interaction, in which the screening is obtained from the direct random-phase approximation (RPA). This amounts to solving the Casida equation
\begin{equation}
\begin{pmatrix}
\mathbf{A} & \mathbf{B} \\
-\mathbf{B} & -\mathbf{A}
\end{pmatrix}
\begin{pmatrix}
\mathbf{X} \\
\mathbf{Y}
\end{pmatrix}
=
\begin{pmatrix}
\mathbf{X} \\
\mathbf{Y}
\end{pmatrix}
\mathbf{\Omega},
\label{eq:casida}
%     \left(\begin{array}{ll}
% \mathrm{A} & \mathrm{~B} \\
% \mathrm{~B} & \mathrm{~A}
% \end{array}\right)\left(\begin{array}{ll}
% \mathrm{X} & \mathrm{Y} \\
% \mathrm{Y} & \mathrm{X}
% \end{array}\right)=\left(\begin{array}{cc}
% \mathrm{X} & -\mathrm{Y} \\
% -\mathrm{Y} & \mathrm{X}
% \end{array}\right)\left(\begin{array}{ll}
% \Omega & 0 \\
% 0 & \Omega
% \end{array}\right)
\end{equation}
where
\begin{align}
    A_{ia,jb} &= \delta_{ij}\delta_{ab} \left(\epsilon_a^{\rm LDA}-\epsilon_i^{\rm LDA}\right) + v(\mathbf{q}),\\
    B_{ia,jb} &= v(\mathbf{q}).
% \left(\green{\epsilon_a^{\rm HF}}+\red{\frac{v_M}{2}}-\green{\epsilon_i^{\rm HF}}-\red{\frac{v_M}{2}}\right)+\langle\Phi_i^a|\green{\hat V_{\rm res}}|\Phi_j^b\rangle\\
%     &= \delta_{ij}\delta_{ab}\left(\green{\epsilon_a^{\rm HF}}-\green{\epsilon_i^{\rm HF}}\right)+\langle\Phi_i^a|\green{\hat V_{\rm res}}|\Phi_j^b\rangle\\
% B_{ia,jb} &= \langle\Phi_i^a|\green{\hat V_{\rm res}}|\Phi_j^b\rangle.
\end{align}
Here, $i$ and $j$ label bands, $a = i + \mathbf{q}$ and $b = j + \mathbf{q}$ label unoccupied orbitals, and $v(\mathbf{q}) = 4\pi/(V q^2)$ is the Coulomb interaction in a supercell of volume $V$.

Because momentum is conserved, the Casida matrix is block diagonal in the transferred momentum $\mathbf{q}$. Diagonalizing each block yields a set of bosonic modes $\nu$ (which include the plasmon) with frequencies $\Omega_\nu(\mathbf{q})$. The corresponding hole--boson coupling is
\begin{equation}
g_{\nu}(\mathbf{k},\mathbf{q}) = \sqrt{2}\sum_{ia} (\mathbf{k}\,,\mathbf{k}-\mathbf{q}\,|\,i\,a)\,(\mathbf{X}+\mathbf{Y})_{ia,\nu}
= \sqrt{2}\,v(\mathbf{q})\sum_{i}(\mathbf{X}+\mathbf{Y})_{i,\,i+\mathbf{q};\,\nu},
\label{eq:coupling}
\end{equation}
where the second equality follows because, for plane waves, every allowed two-electron integral equals $v(\mathbf{q})$; the coupling is therefore independent of $\mathbf{k}$. A naive diagonalization of the full Casida matrix scales as the sixth power of system size. By exploiting momentum conservation, we instead diagonalize a set of much smaller, independent blocks, one per momentum sector $\mathbf{q}$; the cubic symmetry of the electron gas further reduces the number of symmetry-inequivalent sectors that must be treated.

Note that in the electron gas, the LDA exchange--correlation potential is a constant, so the LDA eigenvalues are simply the free-electron kinetic energies shifted by $V^{\rm xc}$. The LDA starting point therefore contains no correlation in its dispersion.
% \subsubsection{TODO: Introduce our RPA theory. }

\subsection{One-shot {GW}+C}
\label{sec:gwc}
In the cumulant formalism, the Green's function is written as
\begin{equation}
G_k(t) = G_k^{0}(t)\, e^{C_k(t)},
\label{eq:cumulant}
\end{equation}
where $G_k^{0}(t)$ is a non-interacting reference Green's function built from $\epsilon_k^{\rm LDA}$, and $C_k(t) = \tilde{C}_k(t) - i\Lambda_k t$ is the cumulant. The static part of the cumulant is $\Lambda_k = \Sigma_k^{\rm x} - V_k^{\rm xc}$, where $\Sigma_k^{\rm x}$ is the exchange self-energy. The reference Green's function and the dynamic part $\tilde{C}_k(t)$ depend on whether one constructs the time-ordered or the retarded cumulant, as described below.

\subsubsection{Time-ordered cumulant}
In the time-ordered cumulant (TC), the dynamic part is
\begin{equation}
    \tilde{C}_k^{\rm T}(t) = \int d\omega\, \frac{\beta_k^{\pm}(\omega)}{\omega^2}\left(e^{-i \omega t}+i \omega t-1\right),
    \label{eq:tc}
\end{equation}
where
\begin{equation}
\beta_k^{\pm}(\omega) = \frac{1}{\pi}\left|\operatorname{Im} \Sigma_k^{\rm R}\!\left(\omega+\epsilon_k^{\rm LDA}\right)\right| \theta\!\left(\pm\left(\omega + \epsilon_k^{\rm LDA} - \epsilon_{\rm F}\right)\right).
\end{equation}
The non-interacting reference is
\begin{equation}
G_k^{{\rm T},0}(t) = \mp i\, \theta(\pm t)\, \theta\!\left(\pm\left(\epsilon_k^{\rm LDA}-\epsilon_{\rm F}\right)\right) e^{-i \epsilon_k^{\rm LDA} t}.
\end{equation}
Throughout, the upper (lower) sign refers to particles with $k>k_{\rm F}$ (holes with $k<k_{\rm F}$), $\epsilon_{\rm F} \equiv \epsilon_{k_{\rm F}}^{\rm LDA}$ is the Fermi energy, and $\theta$ is the Heaviside step function. Thus, for a hole, only the portion of the self-energy below the Fermi level contributes to the cumulant.

\subsubsection{Retarded cumulant}
In the retarded cumulant (RC), the dynamic part is
\begin{equation}
    \tilde{C}_k^{\rm R}(t) = \int d\omega\, \frac{\beta_k^{\rm R}(\omega)}{\omega^2} \left( e^{-i\omega t} + i\omega t - 1 \right),
    \label{eq:rc}
\end{equation}
where
\begin{equation}
\beta_k^{\rm R}(\omega) = \frac{1}{\pi}\left|\operatorname{Im} \Sigma_k^{\rm R}\!\left(\omega+\epsilon_k^{\rm LDA}\right)\right|,
\end{equation}
so that both the particle and hole portions of the self-energy contribute. The non-interacting reference always takes the form
\begin{equation}
G_k^{{\rm R},0}(t) = -i\, \theta(t)\, e^{-i\epsilon_k^{\rm LDA} t}.
\end{equation}
Kas, Rehr, and Reining showed that, in the thermodynamic limit, the retarded cumulant yields a more physical spectral function because it treats particles and holes on an equal footing~\cite{kas2014}. They also showed that, at the bottom of the band, the TC and RC are essentially equivalent.
\subsubsection{BWs2+C}
\subsection{Quasiboson formalism}
The methods discussed below solve the Hamiltonian
\begin{align}
\hat{H}
&=
\hat{H}_{\rm el}
+
\hat{H}_{\rm pl}
+
\hat{H}_{\rm ep}
\end{align}
with $\hat{H}_{\rm el}=  \sum_{\bm{k}\in\mathcal{O}}
\epsilon_{\bm{k}}^{\rm HF}
a_{\bm{k}}^{\dagger}a_{\bm{k}}$, $\hat{H}_{\rm pl}= \sum_{\nu_{\bm{q}}\in\mathcal{M}}
\omega_{\nu_{\bm{q}}}
b_{\nu_{\bm{q}}}^{\dagger}
b_{\nu_{\bm{q}}}$, and $\hat{H}_{\rm ep}= \sum_{\bm{k},\bm{k}+\bm{q}\in\mathcal{O}}
\sum_{\nu_{\bm{q}}\in\mathcal{M}}
g_{\nu}^{\rm pl}(\bm{k},\bm{q})
a_{\bm{k}+\bm{q}}^{\dagger}
a_{\bm{k}}
\left(
b_{\nu_{-\bm{q}}}^{\dagger}
+
b_{\nu_{\bm{q}}}
\right)$, where $\mathcal{O}$ and $\mathcal{M}$ are the sets of occupied bands and plasmonic modes, respectively. As indicated, given a plasmonic mode $\nu$, the associated momentum transfer $\bm{q}$ is already determined. The Plasmonic  frequencies $\omega_{\nu_{\bm{q}}}$and electron-plasmon couplings $g_{\nu}^{\rm pl}(\bm{k},\bm{q})$ are The eigenvalues  and transition densities taken from a direct RPA calculation starting from n. LDA reference as . The seasoned GW practitioner may be wary of this formulation because it seemingly mixes Hartree-fock single-particle energies with LDA screening. However, this has been done with empirical success in the self-consistent GW community for decades, so we pursue the same approach in defining the quasiboson Hamiltonian. 

\subsubsection{Self-consistent cumulant expansion}
For many decades, self-consistency has been pursued has been pursued within the Dyson framework. There have only been a few studies of carrying out the self-consistency in the cumulant formalism, which have shown promise for computing spectral functions of the hosting model . This is the first.  application of the self-consistent cumulant expansion to A fermionic system in the electron gas. Note that this has been done before by Holmes, but this is not true self-consistency, as we show in the appendix.  Specifically, we compute the retarded Green's function in the time domain as In equation X, and then we solve for the cumulant as 
We are always solving
\begin{equation}
\frac{dC_k}{dt}=-i\Lambda_k - \sum_{i\in\mathcal{S}}\sum_\nu|g_{ki\nu}|^2\int_0^t\!d\tau\; e^{-i\Delta_{ki\nu}(t-\tau)}\; e^{-[C_k(t)-C_k(\tau)]+[C_i(t)-C_i(\tau)]},
\end{equation}
where $\Delta_{ki\nu}=\epsilon^{\rm LDA}_i-\epsilon^{\rm LDA}_k-{\Omega_\nu}$, $\Lambda_k$ is a constant, ${g_{ki\nu}}$ is the electron-plasmon coupling, and $\mathcal{S}$ is some set of states.
\subsubsection{Delocalized coupled-cluster theory}
The D2 ansatz is 
We compute the DCC ground state by considering the following processes of excitation out of the reference Hartree-fock reference. Then we apply the tangent-space formalism to compute the spectral functions. We use the nomenclature dCC-X-SY-(Z, A) To denote a method, which, For the ground state calculation is X order In the fermions with Y coupled electron-boson excitations. See this paper for more details. 
\section{Results}
Here, we study finite supercells of the electron gas at rs = 4, converging onto the thermodynamic limit. To begin, we consider the smallest possible supercell with 14 electrons in 19 orbitals (14e19o), for which, We can compare the methods of our study with  the numerically exact result with DMRG, which well coincides wtith CCSDT.\\
For 114e485o, 
\section{Conclusions}


%%%%%%%%%%%%%%%%%%%%%%%%%%%%%%%%%%%%%%%%%%%%%%%%%%%%%%%%%%%%%%%%%%%%%
%% Start the main part of the manuscript here.
%%%%%%%%%%%%%%%%%%%%%%%%%%%%%%%%%%%%%%%%%%%%%%%%%%%%%%%%%%%%%%%%%%%%%




%%%%%%%%%%%%%%%%%%%%%%%%%%%%%%%%%%%%%%%%%%%%%%%%%%%%%%%%%%%%%%%%%%%%%
%% If you are using classical BibTeX rather than biblatex,
%% remove the \printbibliography and uncomment the \bibliograpy one
%%%%%%%%%%%%%%%%%%%%%%%%%%%%%%%%%%%%%%%%%%%%%%%%%%%%%%%%%%%%%%%%%%%%%
\newpage
\bibliography{pqe}


\end{document}
